\documentclass[11pt]{article}

% "review"  -> anónimo, con números de línea (para envío a revisión)
% "final"   -> versión final (nombres de autor, sin números de página)
% "preprint"-> no anónimo, con números de página  <-- usamos esta
\usepackage[preprint]{acl}

% Paquetes estándar
\usepackage{times}
\usepackage{latexsym}
\usepackage[T1]{fontenc}
\usepackage[utf8]{inputenc}
\usepackage{microtype}
\usepackage{inconsolata}
\usepackage{graphicx}

% Español. es-noshorthands evita que babel active atajos (" < > ~ .) que
% chocarían con Tabla~\ref, math, etc.
\usepackage[provide=*,spanish]{babel}
\renewcommand{\abstractname}{Resumen}
\addto\captionsspanish{\renewcommand{\tablename}{Tabla}}

% Para las tablas del trabajo
\usepackage{booktabs}
\usepackage{amsmath}
\usepackage{float}

\title{MiroFish: Simulación Social Generativa y Evaluación de \\ Sistemas Multi-Agente en Escenarios Complejos}

\author{
  \textbf{Andrés Bravo}, \textbf{Joaquín Cadeiras}, \textbf{Bruno De Cruz}, \\
  \textbf{Lucas Ercolano}, \textbf{Eliana Ostrovsky} \\
  Universidad de San Andrés \\
  \small{Procesamiento del Lenguaje Natural}
}

\begin{document}
\maketitle

\begin{abstract}
Los sistemas de simulación social generativa, como MiroFish, pronostican eventos haciendo deliberar a poblaciones de agentes basados en modelos de lenguaje, pero típicamente ejecutan a todos los agentes sobre un mismo modelo. Proponemos poblar la simulación con modelos heterogéneos (Llama-3.3-70B, Gemma-3-27B y Qwen3-8B) bajo la hipótesis de reducir sesgo y costo. Construimos un \textit{benchmark} temporalmente controlado de tres casos —económico, político y deportivo—, lo contrastamos con señales de mercado del mismo corte y sometimos modelos aislados y la red heterogénea a regímenes de estrés. Un modelo aislado y calibrado resultó más preciso y barato; la red heterogénea sin estímulos externos colapsó en ausencia de predicción. La deduplicación del grafo mejoró una corrida, pero no permite generalizar una ganancia predictiva. La heterogeneidad no auto-regula el sesgo y las inyecciones pueden persuadir sin aumentar el acierto. Su valor principal aparece como laboratorio auditable para estudiar dinámicas socio-técnicas, no como reemplazo sistemático de referencias de mercado.
\end{abstract}

\section{Introducción}
\label{sec:intro}

Los grandes modelos de lenguaje (LLM) hicieron viable simular poblaciones de agentes que conversan, se influyen y toman decisiones, abriendo la puerta a emplear estas simulaciones como instrumento de pronóstico de eventos reales \citep{park2023generative, yang2024oasis}. MiroFish \citep{mirofish2025} es un sistema de esta familia: a partir de un corpus semilla instancia una población de agentes en un foro social, los hace deliberar y destila un pronóstico a partir del consenso emergente. La premisa es que, en eventos cuyo desenlace depende de la reacción colectiva, la dinámica social simulada aporta señal más allá de la extrapolación puramente estadística.

Ahora bien, estos entornos suelen ejecutar a todos los agentes sobre un mismo LLM. Una población algorítmicamente homogénea tiende a reproducir las patologías de su único modelo subyacente —cámaras de eco, rigidez de pre-entrenamiento, baja diversidad léxica— y, cuando ese modelo es grande, el costo de cómputo de una simulación completa se vuelve considerable. Postulamos que poblar la simulación con modelos heterogéneos —de distinto tamaño, corpus de pre-entrenamiento y alineación— podría mitigar ambos problemas: la fricción entre arquitecturas divergentes reduciría el sesgo, y la incorporación de modelos más pequeños abarataría la inferencia.

Contrastar esta hipótesis exigía un punto de comparación inexistente: MiroFish no contaba con un \textit{benchmark} que permitiera medir si una modificación mejora o degrada el sistema. Por eso construimos uno propio, con tres casos de distinta temática (económico, político y deportivo) bajo un protocolo estricto de control de fugas de información, y evaluamos primero los modelos en aislamiento (\textit{baseline}) y luego la red multi-modelo, replicando en ambos los mismos regímenes de estrés.

Este trabajo aporta: (i) un mecanismo de ruteo heterogéneo con telemetría y artefactos auditables; (ii) una deduplicación semántica acotada del grafo y una memoria Wiki para el reporte; (iii) un \textit{benchmark} reproducible de tres casos, con \textit{ground truth} aislado y comparación contra señales de mercado temporalmente válidas; y (iv) un análisis empírico de cómo la evidencia, la topología, el modelo y las inyecciones afectan una población homogénea o heterogénea.

\section{Trabajo Relacionado}
\label{sec:related}

\paragraph{Simulación social generativa.} La idea de instanciar agentes basados en LLM que interactúan de forma creíble fue popularizada por los \textit{generative agents} de \citet{park2023generative}, y llevada a escala poblacional por OASIS \citep{yang2024oasis}, que simula hasta un millón de usuarios sobre clones de Twitter y Reddit. MiroFish \citep{mirofish2025} reutiliza este tipo de entorno pero lo orienta al pronóstico de eventos. A diferencia de estos trabajos, que instancian una población homogénea sobre un único modelo, nuestro foco es la heterogeneidad de modelos dentro de una misma simulación y su efecto sobre el sesgo, el costo y la calidad predictiva.

\paragraph{Degeneración y diversidad.} La degradación de sistemas generativos realimentados con su propia salida (\textit{model collapse}) fue caracterizada por \citet{shumailov2024collapse}; observamos un fenómeno análogo cuando la red delibera en circuito cerrado sin entropía exógena. Para cuantificar la variedad lingüística de cada modelo —de forma independiente de su puntería predictiva— empleamos métricas estándar de diversidad: distinct-$n$ \citep{li2016diversity}, Self-BLEU \citep{zhu2018texygen} y Vendi Score \citep{friedman2023vendi}. Por último, dado que las asimetrías de comportamiento entre modelos (proactividad, permeabilidad, cortesía evasiva) se originan en parte en sus esquemas de alineación \citep{ouyang2022instructgpt}, también interpretamos las dinámicas inter-modelo a esa luz.

\section{Metodología}
\label{sec:metodologia}

\subsection{Arquitectura del sistema}

MiroFish \citep{mirofish2025} es un motor de predicción que estima el desenlace de un evento simulando la conversación de una población de agentes en un foro social. A partir de un corpus semilla (noticias, informes, encuestas) el sistema construye un grafo de conocimiento, instancia agentes con perfiles y memoria propios, y los hace interactuar durante un número configurable de rondas; el consenso emergente se destila luego en una predicción. La capa de interacción social reutiliza el entorno OASIS \citep{yang2024oasis}, que clona plataformas tipo Twitter y Reddit e incorpora sistemas de recomendación y un repertorio de acciones (publicar, comentar, seguir, repostear, entre otras).

Sobre esta base introducimos una capa de experimentación y auditoría que permite comparar configuraciones sin perder trazabilidad:

\paragraph{Ruteo y observabilidad.} Por defecto, todos los agentes de una simulación corren sobre un mismo LLM. Implementamos un enrutador que asigna modelos distintos mediante un mapa declarativo en YAML, junto con telemetría para registrar modelo, tokens, latencia, costo y errores. El \textit{Observability Dock} reúne esos datos con los artefactos de cada ejecución, las trazas de \textit{Deep Search} y la memoria Wiki, de modo que una conclusión pueda rastrearse hasta su configuración y evidencia.

\paragraph{Grafo y Deduplicación Semántica.} El grafo base se persiste en Neo4j y conserva entidades, episodios y relaciones. Para mitigar la sobrecarga atencional de los modelos, la etapa posterior a la ingesta emplea \textit{embeddings} para comparar nodos aislados: si la similitud coseno supera un umbral de 0.85, el nodo redundante se descarta.
\paragraph{Memoria a Largo Plazo.} Inspirados en el paradigma del \textit{LLM OS} propuesto por Andrej Karpathy, donde el modelo actúa como CPU y la ventana de contexto como RAM, implementamos una \textbf{Memoria Wiki} que funciona como el sistema de archivos del agente. Esta memoria compila fuentes, afirmaciones, cronología y contradicciones de forma externa y estructurada. Actúa como un artefacto auditable a largo plazo que alimenta el contexto del agente generador de reportes solo cuando es necesario, evitando saturar la memoria de trabajo durante el debate en tiempo real.

Todas las simulaciones se ejecutan en entornos herméticos: los agentes no acceden a Internet y el corpus semilla respeta un corte temporal estricto anterior al evento, evitando fugas de información (\textit{data leakage}). El resultado real se mantiene aislado del input y solo se emplea en la evaluación.

\subsection{Casos de estudio}
\label{subsec:casos}

Al no existir un \textit{benchmark} previo para el sistema, construimos uno propio con tres casos de distinta temática (Tabla~\ref{tab:casos}): económico, político y deportivo. Cada caso plantea un evento cuyo desenlace es posterior al corte documental y verificable de forma objetiva, de modo que el \textit{benchmark} cubre dominios con dinámicas de información heterogéneas manteniendo una evaluación no ambigua. Al representar cada temática con un único evento, el objetivo es la amplitud de dominios y no una medición estadísticamente concluyente dentro de cada uno; una caracterización más fina de las fortalezas y debilidades del sistema por dominio queda como trabajo futuro.

\paragraph{Económico: inflación del IPC (Argentina, 2025).} La variación mensual del Índice de Precios al Consumidor que publica el INDEC. Fijamos el corte documental el 31 de enero de 2025 y pedimos pronosticar el IPC mensual en cuatro horizontes crecientes ($\Delta$1 a $\Delta$4: febrero, abril, julio y diciembre de 2025), para evaluar si el sistema captura la trayectoria de desinflación del programa económico entonces vigente. Es el caso más cuantitativo y cerrado: el desenlace es un valor oficial no ambiguo. El \textit{ground truth} muestra una desinflación real pero no lineal —con un rebote a mitad de año— y una inflación acumulada anual en torno al 31.5\%.

\paragraph{Político: balotaje presidencial (Bolivia, 2025).} La segunda vuelta del 19 de octubre de 2025 entre Rodrigo Paz y Jorge ``Tuto'' Quiroga, en un ciclo marcado por la crisis económica y el colapso electoral del MAS tras casi dos décadas en el poder. A partir de encuestas, resultados de primera vuelta y contexto socioeconómico previos al balotaje, el sistema debe anticipar el ganador y su porcentaje de voto. Es el caso más abierto y polarizado. El resultado real fue el triunfo de Paz con 54.53\% frente al 45.47\% de Quiroga (margen de 9.06 puntos).

\paragraph{Deportivo: final de la Copa América (2024).} La final del 14 de julio de 2024 entre Argentina y Colombia, disputada en Miami. Con un corte documental el día previo (13 de julio), el sistema debe predecir el ganador a partir de previas deportivas, estado de forma e historial. Es un caso simple y auto-verificable, de horizonte muy corto ($\sim$1 día), que opera como prueba de sanidad del protocolo temporal. El desenlace real fue la victoria de Argentina por 1--0 en tiempo suplementario.

\begin{table}[h]
\centering
\small
\resizebox{\columnwidth}{!}{%
\begin{tabular}{@{}llll@{}}
\toprule
\textbf{Tema} & \textbf{Caso} & \textbf{Predicción} & \textbf{Métrica} \\
\midrule
Econ. & IPC Arg. 2025 & Inflación ($\Delta$1--$\Delta$4) & MAE (\%) \\
Polít. & Bolivia 2025 & Ganador / voto & MAE (voto) \\
Dep. & Copa Am. 2024 & Ganador & Acierto \\
\bottomrule
\end{tabular}
}
\caption{Los tres casos del \textit{benchmark}: inflación mensual (INDEC), balotaje presidencial (Paz vs.\ Quiroga) y la final Argentina--Colombia. En todos, el evento es posterior al corte documental y el \textit{ground truth} se mantiene aislado del input.}
\label{tab:casos}
\end{table}

\subsection{Modelos evaluados}

Evaluamos tres LLM de código abierto seleccionados para desacoplar dos variables habitualmente correlacionadas: el tamaño del modelo y la novedad de su fecha de corte. Elegimos deliberadamente una relación invertida entre ambas: \texttt{Llama-3.3-70B} \citep{llama33} es el mayor pero el de corte de conocimiento más antiguo (diciembre de 2023), mientras que \texttt{Qwen3-8B} \citep{qwen3} es el menor y el de corte más reciente (marzo de 2025); \texttt{Gemma-3-27B} \citep{gemma3}, con corte en agosto de 2024, ocupa una posición intermedia en ambos ejes.\footnote{Qwen no publica una fecha de corte oficial para Qwen3; la estimamos en marzo de 2025 sondeando el modelo con una batería de preguntas sobre eventos datados. Debe tomarse como aproximación, no como dato oficial.} Este diseño permite indagar si las patologías que se observen son atribuibles a la escala del modelo o a la cobertura temporal de su pre-entrenamiento.

\subsection{Métricas}

Distinguimos dos familias. Para el \textbf{poder predictivo} (\textit{forecasting}) usamos el Error Absoluto Medio (MAE) respecto del \textit{ground truth}, en la unidad natural de cada caso: puntos porcentuales de inflación (IPC), puntos de voto (Bolivia) y acierto binario del ganador (Copa América). Para el \textbf{comportamiento intrínseco} de los modelos —independiente de su puntería— empleamos métricas agnósticas de diversidad y estabilidad: Self-BLEU y distinct-2 \citep{li2016diversity,zhu2018texygen}, Vendi Score \citep{friedman2023vendi}, entropía categórica y deriva temporal intra-persona, esta última medida con entrevistas paramétricas a los agentes en distintos puntos de la simulación (\textit{Checkpoints Interview}). La capa de telemetría complementa ambas familias registrando el consumo de tokens —nuestra medida de costo de cómputo— y la adherencia al formato de salida (JSON).

% ============================================================================
% RESULTADOS: primero TODO el baseline, después TODO el multimodelo (pedido del
% equipo). Cada "línea" (variable estudiada) es una subsubsección separada.
% ============================================================================
\section{Evaluación}
\label{sec:evaluacion}

\subsection{Baseline: modelos aislados (Single-Agent)}
\label{subsec:baseline}

Como línea de base ejecutamos cada modelo en aislamiento —toda la población sobre un mismo LLM— y sometimos el sistema a una batería de regímenes de estrés. La fase tuvo dos objetivos: medir el poder predictivo en los tres dominios y caracterizar los sesgos de cada arquitectura antes de introducir una red heterogénea. Las pruebas temporales se ejecutaron con \texttt{Gemma-3-27B-IT}; salvo indicación contraria, en las pruebas de rondas y densidad usamos \texttt{Llama-3.3-70B} como modelo ancla.

\subsubsection{Actualización temporal de la evidencia (T0--T3)}

Construimos, para los tres casos, paquetes acumulativos T0--T3 manteniendo fijas 40 rondas y densidad 2. La pregunta, el protocolo y el \textit{ground truth} quedaron fuera del \textit{input}; entre condiciones solo cambió la evidencia disponible hasta cada fecha de corte. Como cada dominio requiere una evaluación distinta, conservamos sus métricas naturales en lugar de reducirlas a un puntaje común.

\paragraph{Bolivia: evidencia electoral, plataformas y encuesta tardía.}
Evaluamos el ganador, los porcentajes de voto, su MAE y el error de margen. El margen es la diferencia en puntos porcentuales entre Paz y Quiroga; su error mide la distancia entre el margen pronosticado y el observado. En T0 el sistema aún debía inferir quiénes llegarían al balotaje: identificó a Luis Arce, presidente saliente asociado al MAS que ya había retirado su candidatura, y a Samuel Doria Medina, candidato opositor que luego quedó eliminado. T1 agregó el resultado de primera vuelta; T2, el contraste entre las plataformas de Paz y Quiroga; y T3, una encuesta tardía favorable a Quiroga junto con evidencia internacional.

\begin{table}[H]
\centering
\small
\begin{tabular}{@{}llll@{}}
\toprule
\textbf{T} & \textbf{Dupla} & \textbf{Ganador} & \textbf{MAE} \\
\midrule
T0 & Arce--Doria & No comparable & -- \\
T1 & Paz--Quiroga & Paz & 2.000 \\
T2 & Paz--Quiroga & Quiroga & 7.020 \\
T3 & Paz--Quiroga & Quiroga & 7.687 \\
\bottomrule
\end{tabular}
\caption{Bolivia por corte. T0 no identifica el balotaje real; la evidencia completa se detalla en la Tabla~\ref{tab:app_seed_bolivia}.}
\label{tab:seed_bolivia}
\end{table}

T0 no falló solo al elegir un ganador: construyó un balotaje imposible y, por ello, no tiene MAE de voto ni error de margen evaluables. El resultado de primera vuelta corrigió el campo de candidatos y T1 produjo a Paz como ganador. T2 y T3 volvieron a Quiroga; la encuesta tardía tuvo más peso que la evidencia anterior. Acercar el corte al desenlace, por tanto, no mejoró necesariamente el pronóstico.

\paragraph{IPC: información mensual, expectativas y política económica.}
Comparamos las proyecciones mensuales de febrero, abril, julio y diciembre, además del acumulado anual de 2025. Cuando el informe devolvió un rango, usamos su punto medio para calcular el error absoluto y el MAE. T0 contiene el contexto social, institucional y macroeconómico de 2024; T1 suma el REM del BCRA y la evaluación \textit{ex post} del FMI; T2 incorpora el IPC oficial de diciembre de 2024; y T3 agrega el cambio del \textit{crawling peg} al 1\%, el superávit fiscal y nuevas señales monetarias.

\begin{table}[H]
\centering
\small
\begin{tabular}{@{}lcc@{}}
\toprule
\textbf{T} & \textbf{MAE mensual} & \textbf{MAE con acumulado} \\
\midrule
T0 & 1.538 & 4.03 \\
T1 & 1.250 & 2.30 \\
T2 & 1.200 & 2.26 \\
T3 & \textbf{0.900} & \textbf{1.52} \\
\bottomrule
\end{tabular}
\caption{IPC por corte; composición documental en la Tabla~\ref{tab:app_seed_ipc}.}
\label{tab:seed_ipc}
\end{table}

En esta serie la reducción fue sostenida: el MAE mensual bajó de 1.538 en T0 a 0.900 en T3, y el MAE que incorpora el acumulado anual cayó de 4.03 a 1.52. Los datos oficiales anclaron el corto plazo, mientras que las expectativas y las señales cambiarias, fiscales y monetarias ayudaron a ajustar el sendero posterior.

\paragraph{Copa América: forma reciente, mercado y contrapesos.}
Evaluamos por separado el campeón, la confianza declarada, la probabilidad puntual, su intervalo y el margen esperado de gol. Un margen de 1 representa una victoria por un gol. T0 contiene antecedentes, planteles y contexto; T1 suma semifinales y forma reciente; T2 incorpora cuotas, modelos y previas favorables a Argentina; y T3 agrega contrapesos sobre Colombia y evidencia táctica.

El campeón y el margen esperado de un gol permanecieron idénticos en T0--T3: Argentina por un gol, en línea con el 1--0 observado. Por esa invariancia no repetimos una tabla en el cuerpo. La confianza declarada osciló entre 0.72 y 0.78, mientras que la probabilidad puntual pasó de 0.625--0.650 a 0.490; la evidencia tardía recalibró la incertidumbre sin cambiar la conclusión. La salida completa se conserva en la Tabla~\ref{tab:app_seed_copa}.

En conjunto, los paquetes temporales no dibujan una curva universal de mejora. En IPC, las nuevas observaciones redujeron el error de manera acumulativa; en Bolivia, una encuesta tardía desplazó el acierto de T1; y en Copa América, la predicción se mantuvo mientras se ajustaba la incertidumbre. El \textit{seed data} debe evaluarse por la naturaleza y el momento de la evidencia, no solo por su cercanía al desenlace.

\subsubsection{Escalamiento topológico (rondas y densidad)}

Variamos la profundidad del debate ($R$) y la densidad de interacción ($D$), manteniendo fijo el resto del protocolo. El efecto dependió del dominio y no apareció una configuración universal.

\paragraph{Bolivia.}
Aumentar las rondas no corrigió el pronóstico hacia Quiroga. R40-D1 incluso elevó el MAE a 10.000, mientras que R10-D2 y R80-D2 produjeron el mismo MAE de 7.687.

\paragraph{IPC.}
En esta serie, en cambio, más rondas mejoraron la proyección: el MAE pasó de 3.12\% en R10-D2 a 1.84\% en R80-D2. Con 40 rondas, D2 superó a D1 y D3, aunque el costo y la latencia crecieron. Las tres réplicas de R80-D2 obtuvieron 1.84\%, 1.91\% y 1.80\% (media 1.85\%, $\sigma=0.05$).

\paragraph{Copa América.}
Las cinco condiciones mantuvieron exactamente la misma salida: Argentina, probabilidad puntual 0.475 y margen esperado de un gol. La evidencia disponible dominó el efecto topológico: más deliberación no produjo una ganancia observable. Las matrices completas de Bolivia e IPC se trasladan al Apéndice~\ref{app:extended_tables}.

\subsubsection{Sesgos por arquitectura}

Fijada una configuración de control ($R=40$, $D=2$), comparamos los tres modelos en el caso cuantitativo. La precisión ordena los modelos por tamaño: \texttt{Llama-3.3-70B} obtiene el menor error (MAE 2.31\%), seguido de \texttt{Gemma-3-27B} (2.67\%) y \texttt{Qwen3-8B} (3.45\%). En su mejor configuración ($R=80$, $D=2$), Llama es además estable entre réplicas (MAE medio 1.85\%, $\sigma = 0.05$). Sin embargo, en escenarios sociales abiertos cada arquitectura exhibe una patología distinta: Llama tiende a la confirmación recursiva (burbuja epistemológica); Gemma muestra una impermeabilidad factual —mantuvo su voto pese a las estadísticas en contra en el caso deportivo, por el peso de su pre-entrenamiento—; y Qwen cae en una pasividad neutral, aunque es el más permeable a inyecciones de datos. Estas asimetrías motivan directamente el diseño multi-modelo (§\ref{subsec:multimodelo}).





\subsubsection{Benchmark frente a señales de mercado}

Para evitar que la evaluación se limite a comparar MiroFish contra el desenlace, alineamos cada salida T0--T3 con una señal externa disponible en el mismo corte: encuestas en Bolivia, REM y proxies de mercado para IPC, y probabilidades implícitas de apuestas para Copa. No ejecutamos nuevas simulaciones; el análisis reevalúa los artefactos ya guardados. De 28 filas, 27 fueron comparables y MiroFish obtuvo menor error normalizado en 9 (Tabla~\ref{tab:market_baseline}); el delta Brier medio fue $+0.014$, donde un valor positivo favorece al proxy.

\begin{table}[H]
\centering
\small
\begin{tabular}{@{}lccp{3.0cm}@{}}
\toprule
\textbf{Caso} & \textbf{Comp.} & \textbf{Gana MF} & \textbf{Lectura} \\
\midrule
Bolivia & 3/4 & 1 & Ambos siguen la señal tardía errónea \\
IPC & 20/20 & 7 & Proxy mejor temprano; MF mejora en T3 \\
Copa & 4/4 & 1 & MF acierta, pero queda subcalibrado \\
\midrule
Total & 27/28 & 9 & $\overline{\Delta\text{Brier}}=+0.014$ \\
\bottomrule
\end{tabular}
\caption{MiroFish frente a señales externas del mismo corte. ``Gana MF'' cuenta filas con menor error normalizado; las proxies difieren en calidad y no constituyen un mercado homogéneo.}
\label{tab:market_baseline}
\end{table}

La comparación aporta contexto, no un veredicto único. En Bolivia, tanto MiroFish como la encuesta tardía favorecieron a Quiroga aunque ganó Paz; en IPC, las señales de mercado dominaron al inicio, pero MiroFish redujo el error en varias proyecciones T3; y en Copa, el sistema mantuvo al campeón correcto pero asignó menos probabilidad que el mercado. En conjunto, el modelo no supera de forma sistemática a referencias simples y su valor debe medirse contra ellas, no solo contra acierto/fallo final.

\subsubsection{Sensibilidad a Inyecciones Programadas y Ruido}

Extendimos la inyección programada a Copa América, Bolivia e IPC para evaluar la vulnerabilidad de los modelos ante estímulos externos, separando contenido, momento y cercanía temática. Cruzamos Gemma y Llama con control; señal favorable temprana, media y tardía; contraseñal media; ruido cercano; y ruido ajeno. Las 42 corridas fueron válidas: los seis controles dispararon cero eventos y las 36 condiciones inyectadas, exactamente uno (Tabla~\ref{tab:s3_cross_topic}).

Un hallazgo crítico sobre esta vulnerabilidad ocurrió en el dominio económico bajo alta profundidad (R = 80, D = 2). Al introducir ruido financiero—información no predictiva— el modelo asimiló rápidamente la especulación (\textit{herd behavior}). Aunque en esa variante aislada alcanzó un MAE excepcional de 0.9875\% al anclar un escenario recesivo, esta coincidencia no convierte al ruido en evidencia útil: demuestra empíricamente la susceptibilidad del sistema al momento de inyección y a la saliencia del estímulo.
\begin{table}[H]
\centering
\scriptsize
\begin{tabular}{@{}lp{2.0cm}p{3.5cm}@{}}
\toprule
\textbf{Caso} & \textbf{Control} & \textbf{Mayor sensibilidad} \\
\midrule
Copa & Argentina & Contraseñal: Llama vira débilmente a Colombia \\
Bolivia & Indef./Paz & Contraseñal: ambos a Quiroga; ruido afecta a Llama \\
IPC & Desinflación & Contraseñal: Gemma indef.; Llama proyecta rebote \\
\bottomrule
\end{tabular}
\caption{Resumen de 42 corridas S3. La matriz completa está en la Tabla~\ref{tab:app_s3_cross_topic}.}
\label{tab:s3_cross_topic}
\end{table}

La contraseñal produjo el cambio más consistente: Llama viró a Colombia, ambos modelos a Quiroga y, en IPC, Gemma quedó indefinido mientras Llama proyectó un rebote. Gemma permaneció rígido en fútbol. El momento de la señal favorable cambió las menciones, no su categoría final.

El ruido dependió del dominio: no alteró Copa ni IPC, pero Bolivia fue vulnerable. Con Llama, el ruido político cercano desplazó Paz hacia Quiroga y el ajeno eliminó una dirección clara. La robustez, por tanto, no se generaliza entre casos ni modelos.

\subsubsection{Diversidad léxica y estabilidad (evaluación intrínseca)}

Con independencia de la puntería predictiva, medimos el comportamiento generativo intrínseco de cada modelo con métricas agnósticas, sobre el caso más abierto (Bolivia) y con control de varianza (tres corridas por modelo con \textit{equal-N}). Qwen3 no pudo incluirse por falta de acceso al modelo, de modo que la comparación es entre Gemma y Llama (Tabla~\ref{tab:intrinseca}).

Surgen tres asimetrías robustas, con rangos no solapados entre modelos. \textbf{Proactividad}: Llama es marcadamente reactivo —57 comentarios por cada 10 publicaciones originales—, mientras que Gemma es proactivo (18 publicaciones, 16 comentarios). \textbf{Diversidad léxica}: los posts de Gemma son mucho más variados (distinct-2 0.87 frente a 0.33; Self-BLEU 0.11 frente a 0.73), en tanto que Llama recicla estructuras. \textbf{Deriva temporal}: las personas de Gemma evolucionan a lo largo de la simulación (distancia \textit{embedding} inicio-fin de 0.41), mientras que las de Llama permanecen estáticas (0.09).

Es notable que estas asimetrías no se explican por las personas iniciales: las de Llama son incluso algo más diversas al crearse (Vendi 9.2 frente a 8.4), y aun así sus posts convergen y se repiten. La menor diversidad y la estabilidad de Llama son, por tanto, un fenómeno de generación y no una consecuencia de personas homogéneas.

\begin{table}[t]
\centering
\small
\begin{tabular}{@{}lccc@{}}
\toprule
\textbf{Modelo} & \textbf{distinct-2}$\uparrow$ & \textbf{Self-BLEU}$\downarrow$ & \textbf{Deriva}$\uparrow$ \\
\midrule
Gemma-3-27B & \textbf{0.87} & \textbf{0.11} & 0.41 \\
Llama-3.3-70B & 0.33 & 0.73 & 0.09 \\
\bottomrule
\end{tabular}
\caption{Evaluación intrínseca de la salida (Caso Bolivia; \textit{variance-checked}, 3 corridas/modelo, \textit{equal-N}). Gemma produce texto más diverso y sus personas derivan más en el tiempo; Llama es más repetitivo y estático. Qwen3 no pudo evaluarse. Deriva: distancia \textit{embedding} intra-persona inicio-fin ($\uparrow$ = más deriva).}
\label{tab:intrinseca}
\end{table}

\subsubsection{Síntesis del baseline}

La fase de aislamiento mostró que ninguna variable mejora todos los dominios por igual. La recencia de la evidencia ayudó al IPC, corrigió transitoriamente Bolivia y solo recalibró la incertidumbre deportiva; más rondas mejoraron el caso cuantitativo, pero no corrigieron el sesgo electoral; y las inyecciones S3 revelaron sensibilidad específica al caso, al modelo y a la cercanía temática. El ruido no constituyó una fuente confiable de señal. Estas asimetrías motivaron evaluar si una población heterogénea podía compensar las limitaciones de cada modelo.

\subsection{Estudios de Ablación de Ingesta y Grafo}

\subsubsection{Deduplicación de nodos: costo y precisión}

El grafo real del caso Copa contiene 26 entidades y 107 relaciones; incluso allí aparecen variantes nominales como \textit{Copa America 2024}/\textit{Copa América 2024} y dos nodos \textit{Opta}. La optimización implementada actúa solo sobre nodos aislados: los compara con entidades conectadas y descarta el candidato si la similitud coseno supera 0.85. En una corrida del caso IPC, la variante con ingesta manual y deduplicación redujo el MAE de 2.31\% a 1.475\%. Es una señal favorable de una única corrida, no evidencia de una mejora predictiva general ni una medición exacta de reducción del grafo.


\subsubsection{Permeabilidad e ingesta autónoma (Deep Search)}

Evaluamos una ingesta totalmente autónoma (\textit{Deep Search}) con corte temporal estricto. En la corrida IPC, el MAE pasó de 2.31\% con ingesta manual a 2.475\% con búsqueda pura: solo se recuperaron dos entidades poco relevantes y la proyección perdió anclas documentales. El resultado no invalida la búsqueda como herramienta de apoyo, pero favorece un diseño híbrido: corpus curado como base, búsqueda trazable para cubrir vacíos y revisión mediante Observability/Wiki.


\subsection{Multi-modelo: red heterogénea}
\label{subsec:multimodelo}

Motivados por las asimetrías del baseline, poblamos una misma simulación con los tres modelos y replicamos los mismos regímenes de estrés. Nuestra hipótesis era que la fricción entre arquitecturas divergentes funcionaría como auto-regulación y reduciría el sesgo.

\subsubsection{Escalamiento de rondas y colapso semántico}

En el baseline, escalar las rondas con un solo modelo (Llama) inflaba la cámara de eco hasta un MAE de 10.0 en el balotaje de Bolivia. Al replicar la prueba con la red heterogénea y sin estímulos externos, la hipótesis de auto-regulación se refutó de forma contundente: en lugar de corregir hacia la verdad, el sistema sufrió un \textbf{colapso semántico} (\textit{model collapse}). Forzados a dialogar sin nueva entropía, los agentes cayeron en un \textit{deadlock} dialéctico —la hiper-reactividad de Llama frente a la proactividad de Gemma, moduladas por sus esquemas de alineación— y el foro se llenó de cortesía vacía (p.\ ej.\ \textit{``I completely agree with the importance of considering the candidates' policies...''}), sin que ningún modelo tomara postura.

El efecto es más severo que una simple pérdida de precisión: el 100\% de los agentes evadió emitir un voto, de modo que el evaluador no pudo extraer ningún pronóstico válido. La métrica de error resultante (66.667) no es, por tanto, un MAE de pronóstico comparable con los del baseline, sino el valor mecánico que produce el evaluador ante una predicción nula —el 100\% de la masa asignada a la categoría ``otros''— y la registramos como inválida (Tabla~\ref{tab:collapse}). El hallazgo cualitativo es que, en cautiverio y sin entropía externa, la red heterogénea no colapsa hacia el error sino hacia la \textit{ausencia de predicción}.

\begin{table}[t]
\centering
\small
\begin{tabular}{@{}llc@{}}
\toprule
\textbf{Arquitectura} & \textbf{Predicción} & \textbf{MAE} \\
\midrule
Baseline (Llama), R10 & Quiroga & 7.687 \\
Baseline (Llama), R40 & Quiroga & 10.000 \\
Multi-modelo, R10 & Sin predicción & 66.667$^{*}$ \\
Multi-modelo, R40 & Sin predicción & 66.667$^{*}$ \\
\bottomrule
\end{tabular}
\caption{Balotaje de Bolivia sin inyecciones externas. La red heterogénea colapsa en abstención total: no emite un pronóstico válido. $^{*}$Valor mecánico de una predicción nula (100\% ``otros''), marcado como inválido en el informe fuente; no es comparable con los MAE del baseline.}
\label{tab:collapse}
\end{table}

\paragraph{Profundidad y costo en IPC trimodelo.} La telemetría muestra que escalar la red tampoco produjo una mejora monotónica: R10 obtuvo 3/5 con 50\,231 tokens, mientras R80 cayó a 1/5 con 5\,546\,655 tokens y acumuló 11 errores y 39 fallas de parseo. Más deliberación multiplicó el costo en más de dos órdenes de magnitud sin aumentar la calidad evaluada.

\subsubsection{Continuación de la inyección de ruido: shock tardío}

Como continuación de la prueba de ruido del baseline, y sabiendo que la red heterogénea colapsaba en neutralidad, inyectamos una encuesta falsa tardía (\textit{Late Quiroga Lead Poll}) a mitad del debate inactivo. El shock desestabilizó el estado de cortesía: Qwen, con la permeabilidad observada en aislamiento, absorbió el estímulo de inmediato, lo tradujo en retórica agresiva e impulsó la polarización del foro. En esta corrida, la señal externa reactivó la deliberación, pero ese cambio de dinámica no demuestra que el resultado se haya vuelto más verdadero o robusto.

\section{Conclusiones}
\label{sec:conclusiones}
El contraste entre el \textit{baseline} aislado y la red heterogénea nos permite volver sobre las dos hipótesis que motivaron este trabajo —que la heterogeneidad de modelos reduciría el sesgo y el costo— y matizarlas a la luz de la evidencia.

\paragraph{Para \textit{forecasting} puro, la heterogeneidad no compensa.} Un modelo aislado y bien calibrado alcanzó MAE 1.84\% en IPC y fue más barato. En la red trimodelo, pasar de R10 a R80 elevó el consumo de 50\,231 a 5\,546\,655 tokens mientras el puntaje cayó de 3/5 a 1/5; en Bolivia, la red sin estímulos no emitió un pronóstico utilizable (§\ref{subsec:multimodelo}). La deduplicación fue una optimización ortogonal: en una corrida redujo el MAE a 1.475\%, pero no medimos de forma independiente su ahorro ni podemos generalizar el efecto.

\paragraph{La diversidad de modelos no se auto-regula, pero es permeable a la intervención.} En circuito cerrado, la red no converge a la verdad sino a un \textit{deadlock} de cortesía. Una encuesta falsa tardía reactivó la deliberación, pero no demostró mayor veracidad. La matriz S3 confirma la distinción: una contraseñal cambió direcciones en los tres dominios, mientras que el ruido tuvo efectos específicos al modelo y al caso. Permeabilidad no equivale a robustez ni acierto.

\paragraph{MiroFish debe competir contra referencias externas.} En 27 comparaciones temporales, MiroFish superó al proxy en 9 y el delta Brier medio favoreció levemente a las señales de mercado. El resultado depende de la calidad del proxy, pero evita celebrar un acierto que una referencia simple ya anticipaba y vuelve explícita la calibración pendiente.

\paragraph{La precisión predictiva no requiere diversidad generativa.} Un hallazgo transversal es que el mejor pronosticador (Llama) es, a la vez, el modelo intrínsecamente menos diverso y más estable en el tiempo, mientras que el más diverso (Gemma) es también el más volátil, con personas que derivan más a lo largo de la simulación (§\ref{subsec:baseline}). Puntería y riqueza generativa son ejes independientes, y optimizar uno no implica el otro.

\paragraph{El valor de la red heterogénea es como laboratorio socio-técnico.} El ruteo multi-modelo, la telemetría, la Wiki y el panel de observabilidad permiten estudiar polarización, absorción de (des)información y ruptura de consensos con trazabilidad de evidencia, modelo y costo. Ese montaje es más convincente como instrumento experimental que como pronosticador autónomo.

% La sección de Limitaciones es OBLIGATORIA en ACL y NO cuenta para el límite
% de 8 páginas.
\section*{Limitaciones}

Este trabajo tiene un alcance exploratorio y sus conclusiones deben leerse con varias salvedades.

\paragraph{Cobertura muestral.} El \textit{benchmark} usa un único evento por temática (económico, político y deportivo). Esto permite abarcar dominios diversos, pero impide afirmaciones estadísticamente concluyentes sobre las fortalezas o debilidades del sistema en cada uno; distinguir un efecto de dominio del ruido de un caso particular requeriría varios eventos por temática.

\paragraph{Control de varianza parcial.} Solo algunas mediciones se repitieron con réplicas y control de varianza (la precisión del IPC y el análisis intrínseco de diversidad y deriva). S3 cubre 42 celdas, pero cada combinación de caso, modelo y condición tiene una sola corrida. El ruido financiero en alta profundidad y el colapso multi-modelo también provienen de corridas únicas, por lo que sus valores puntuales deben tomarse como indicativos y no como estimaciones robustas.

\paragraph{Cobertura de modelos.} La comparación intrínseca y la matriz S3 se realizaron entre Gemma-3 y Llama-3.3; Qwen3 no pudo incluirse en esas fases. Además, evaluamos una única terna de modelos en un rango de tamaño acotado (8B--70B) y una sola configuración de mezcla. No podemos garantizar que las conclusiones sean agnósticas al tamaño ni a la elección concreta de modelos; explorar otras combinaciones y escalas mayores es trabajo pendiente.

\paragraph{Métricas de evaluación.} El MAE de participación electoral se calcula a partir de la extracción, mediante expresiones regulares, de porcentajes en el informe generado; una predicción mal formada puede no ser detectada. La matriz S3 usa conteos determinísticos de términos en publicaciones y comentarios, que incluyen el documento inyectado y no equivalen a un juicio semántico del ReportAgent. En particular, el valor asociado al colapso semántico multi-modelo no es un error de pronóstico comparable, sino el resultado mecánico de una predicción nula, y así lo señalamos. El MAE se expresa además en unidades distintas según el caso, por lo que los valores no son comparables entre dominios.

\paragraph{Baseline de mercado.} Las referencias externas no tienen calidad uniforme: Bolivia combina señales de primera vuelta y encuestas de balotaje; IPC usa REM y promedios implícitos para algunos meses; y Copa incluye una proxy pre-torneo en T0. El delta Brier agregado usa la normalización definida por el benchmark y debe leerse junto con el desglose por caso, no como una prueba estadística de superioridad.

\paragraph{Reproducibilidad.} Las simulaciones dependen de modelos servidos por terceros y el costo se aproxima por el conteo de tokens de la telemetría, no por un perfilado exhaustivo ni por su traducción a un costo monetario. La deduplicación de nodos, por su parte, se validó en un único caso, sin cuantificar de forma independiente la reducción exacta del tamaño del grafo.

\bibliography{custom}

\clearpage
\onecolumn
\appendix
\section{Tablas extendidas}
\label{app:extended_tables}

Las tablas del cuerpo priorizan la comparación directa. Aquí conservamos la composición documental y las matrices completas empleadas para interpretar cada corrida.

\begin{table}[H]
\centering
\small
\begin{tabular}{@{}lp{4.5cm}p{4.9cm}p{1.5cm}p{1.6cm}@{}}
\toprule
\textbf{T} & \textbf{Evidencia añadida} & \textbf{Candidatos y ganador} & \textbf{MAE voto} & \textbf{Error margen} \\
\midrule
T0 & Contexto y encuestas previas a la primera vuelta & Arce (fuera de la boleta) y Doria Medina; no identifica el balotaje real & -- & -- \\
T1 & Primera vuelta: Paz primero, Quiroga segundo y MAS fuera & Paz y Quiroga; gana Paz & 2.000 & 0.06 \\
T2 & Plataformas: Paz moderado y Quiroga pro-mercado/FMI & Paz y Quiroga; gana Quiroga & 7.020 & 18.06 \\
T3 & Encuesta tardía favorable a Quiroga y contexto internacional & Paz y Quiroga; gana Quiroga & 7.687 & 18.06 \\
\bottomrule
\end{tabular}
\caption{Bolivia T0--T3. T0 no admite métricas comparables porque no identifica la dupla real Paz--Quiroga.}
\label{tab:app_seed_bolivia}
\end{table}

\begin{table}[H]
\centering
\small
\begin{tabular}{@{}lp{9.4cm}cc@{}}
\toprule
\textbf{T} & \textbf{Evidencia añadida} & \textbf{MAE mensual} & \textbf{MAE con acumulado} \\
\midrule
T0 & Base social, institucional y macroeconómica de 2024 & 1.538 & 4.03 \\
T1 & REM del BCRA, evaluación \textit{ex post} del FMI y expectativas & 1.250 & 2.30 \\
T2 & IPC oficial de diciembre de 2024 y marco político de 2025 & 1.200 & 2.26 \\
T3 & \textit{Crawling peg}, superávit fiscal e informes BCRA/Banco Mundial & 0.900 & 1.52 \\
\bottomrule
\end{tabular}
\caption{IPC T0--T3: composición de cada paquete y errores.}
\label{tab:app_seed_ipc}
\end{table}

\begin{table}[H]
\centering
\small
\resizebox{\textwidth}{!}{%
\begin{tabular}{@{}lp{5.0cm}ccccc@{}}
\toprule
\textbf{T} & \textbf{Evidencia añadida} & \textbf{Campeón} & \textbf{Conf.} & \textbf{Prob.} & \textbf{Intervalo} & \textbf{Margen} \\
\midrule
T0 & Antecedentes, planteles y contexto de la final & Argentina & 0.75 & 0.625 & [0.575, 0.675] & 1 [0.5, 1.5] \\
T1 & Semifinales y forma reciente & Argentina & 0.75 & 0.650 & [0.600, 0.700] & 1 [0.5, 1.5] \\
T2 & Cuotas, modelos y previas favorables a Argentina & Argentina & 0.78 & 0.490 & [0.460, 0.510] & 1 [0, 2] \\
T3 & Contrapesos de Colombia y evidencia táctica & Argentina & 0.72 & 0.490 & [0.460, 0.510] & 1 [1, 2] \\
\bottomrule
\end{tabular}%
}
\caption{Copa América T0--T3; la final observada fue Argentina 1--0 Colombia.}
\label{tab:app_seed_copa}
\end{table}

\begin{table}[H]
\centering
\begin{minipage}[t]{0.46\textwidth}
\centering
\small
\begin{tabular}{@{}lcc@{}}
\toprule
\textbf{Condición} & \textbf{Ganador} & \textbf{MAE voto} \\
\midrule
R10-D2 & Quiroga & 7.687 \\
R40-D1 & Quiroga & 10.000 \\
R40-D2 & Quiroga & 7.820 \\
R40-D3 & Quiroga & 7.820 \\
R80-D2 & Quiroga & 7.687 \\
\bottomrule
\end{tabular}
\end{minipage}\hfill
\begin{minipage}[t]{0.50\textwidth}
\centering
\small
\begin{tabular}{@{}lccc@{}}
\toprule
\textbf{Cond.} & \textbf{MAE} & \textbf{Costo} & \textbf{Lat.} \\
\midrule
R10-D2 & 3.12\% & USD\,0.0214 & 247\,s \\
R40-D1 & 2.45\% & USD\,0.0856 & 338\,s \\
R40-D2 & 2.31\% & USD\,0.0856 & 675\,s \\
R40-D3 & 2.89\% & USD\,0.0856 & 309\,s \\
R80-D2 & 1.84\% & USD\,0.1712 & 3277\,s \\
\bottomrule
\end{tabular}
\end{minipage}
\caption{Matrices completas de rondas y densidad: Bolivia (izquierda) e IPC (derecha).}
\label{tab:app_rounds}
\end{table}

\begin{table}[H]
\centering
\scriptsize
\begin{tabular}{@{}p{2.2cm}p{1.9cm}p{2.4cm}p{2.4cm}p{2.4cm}p{2.3cm}@{}}
\toprule
\textbf{Caso / modelo} & \textbf{Control} & \textbf{Señal favorable} & \textbf{Contraseñal media} & \textbf{Ruido cercano} & \textbf{Ruido ajeno} \\
\midrule
Copa / Gemma & Argentina & Argentina & Argentina (se estrecha) & Argentina & Argentina \\
Copa / Llama & Argentina & Argentina & Colombia (cambio débil) & Argentina & Argentina \\
Bolivia / Gemma & Indefinido & Paz & Quiroga & Indefinido & Indefinido \\
Bolivia / Llama & Paz & Paz & Quiroga & Quiroga & Indefinido \\
IPC / Gemma & Desinflación & Desinflación & Indefinido & Desinflación & Desinflación \\
IPC / Llama & Desinflación & Desinflación & Rebote inflacionario & Desinflación & Desinflación \\
\bottomrule
\end{tabular}
\caption{Dirección heurística de las 42 corridas S3 (3 casos $\times$ 2 modelos $\times$ 7 condiciones). ``Señal favorable'' resume los tres momentos, que conservaron la misma dirección categórica. La métrica cuenta menciones en la conversación, incluido el documento inyectado; audita presión y contaminación, pero no reemplaza una evaluación semántica del informe final.}
\label{tab:app_s3_cross_topic}
\end{table}

\end{document}
